\section{扩展功能}

\subsection{方向参数}

方向参数控制电机转动方向的引脚电平。

\begin{table}[H]
    \centering
    \caption{方向参数默认值}
    \label{tab:direction-default}
    \begin{tabular}{lc}
        \toprule
        \textbf{参数} & \textbf{默认值} \\
        \midrule
        电机 A & 1 \\
        电机 B & 0 \\
        \bottomrule
    \end{tabular}
\end{table}

\textbf{一般情况下不要修改方向参数。}只有在现场发现电机实际转动方向和预期相反时，才考虑调整。通常是接线或机械安装的问题，优先排查硬件；确认需要改方向时，在下拉框选 \texttt{0} 或 \texttt{1}，然后点"保存A"或"保存B"。

保存后写入设备，断电保留。

\begin{figure}[H]
    \centering
    \includegraphics[width=0.7\textwidth]{figures/ui-direction.png}
    \caption{方向参数区域。电机 A 和电机 B 的方向设置与保存按钮。}
    \label{fig:ui-direction}
\end{figure}

\subsection{启动信号读取}

点击"读取启动信号"，可以查看外部启动输入脚当前的电平状态。

\begin{table}[H]
    \centering
    \caption{启动信号状态说明}
    \label{tab:signal-status}
    \begin{tabular}{lp{8cm}}
        \toprule
        \textbf{显示} & \textbf{含义} \\
        \midrule
        HIGH & 外部启动输入为高电平，设备可被外部信号触发 \\
        LOW & 外部启动输入为低电平，没有检测到外部启动输入 \\
        \bottomrule
    \end{tabular}
\end{table}

如果外部启动没有反应，可以用这个按钮排查：先点一次，记下显示的状态；然后触发外部启动源，再点一次，对比两次结果。如果两次都是 LOW，说明信号没有送到控制板，检查外部信号线。如果第二次变成了 HIGH，但设备仍不运行，反馈研发排查固件状态和电机驱动。
